\section{An intrinsic no-loss boundary class and its mixed top calculus}
\label{sec:r11-capacity-tail-boundary-class}

Retain the original field, analytic space, arbitrary background $S$,
ordinary domain $V$, and $q=n-p\geq1$ from
Section~\ref{sec:r11-unbounded-comparison}. Every potential below is
negative extended classical psh, with its actual values retained.
The capacity always has the fixed tester domain $V$. The results
of this section concern top-degree measures. They do not construct
all lower-degree products without the separate energy hypotheses.

\begin{lemma}[Compact conservation of truncation mass]
\label{lem:r11-compact-truncation-conservation}
If $\mathrm{BG}(u,0)$ holds, then $\mu_u$ in (UC3) is locally
finite, without assuming absence of truncation loss. More precisely,
fix $M>0$ and $H\Subset V$, set
$K_M=\overline{\{u<-M\}}\Subset V$, and choose an original smooth
compact function $0\leq\chi\leq1$ which is one on an ordinary
neighborhood of $H\cup K_M$. Then for integers $j,l>M$,
\[
                       \int\chi\,d\mu_{u,j}
                           =\int\chi\,d\mu_{u,l}.
 \tag{BC1}
\]
In particular $\sup_j\mu_{u,j}(H)<\infty$.
\end{lemma}

\noindent\textbf{Status: \texttt{NEW\_PROVED}} as a compact
conservation statement for unbounded actual potentials, using only
the frozen bounded telescoping identity and original compact Stokes.

\begin{proof}
Outside $K_M$ one has $u\geq-M$, hence $u_j=u_l=u$ there.
The bounded mixed calculus gives
\[
 \mu_{u,j}-\mu_{u,l}
   =D\left((u_j-u_l)
       \sum_{r=0}^{q-1}(Du_j)^r(Du_l)^{q-1-r}\wedge S\right).
 \tag{BC2}
\]
Every product in the parentheses is a frozen original positive
closed product. Its bounded Borel multiplier is zero outside
$K_M$, so the entire current in parentheses has support in $K_M$.
Pair (BC2) with $\chi$. Since $D\chi$ vanishes near $K_M$,
the pairing is zero, proving (BC1). Fixing one $l>M$ bounds
the nonnegative masses on $H$ for every $j>M$; the remaining
finitely many $j$ have finite mass. The increasing restricted
measures in (UC3) have the same bound on $H$, proving local
finiteness of their limit. No bound for a lower-degree truncation
product is inferred from this top-degree Stokes argument.
\end{proof}

\begin{definition}[The no-loss zero-boundary class]
\label{def:r11-no-loss-boundary-class}
Set
\[
 \mathcal F_{S,\partial}(V)
   =\{u\in\mathcal F_{S,\mathrm{loc}}(V):
                                      \mathrm{BG}(u,0)\}.
 \tag{BC3}
\]
The notation records compact strict negative sublevels, not an
unstated ordinary boundary regularity hypothesis.
\end{definition}

\begin{theorem}[An exact capacity-tail criterion]
\label{thm:r11-capacity-tail-criterion}
Let $u$ satisfy $\mathrm{BG}(u,0)$. The following are equivalent:
\begin{enumerate}
\item $u\in\mathcal F_{S,\partial}(V)$;
\item $\displaystyle\lim_{t\to\infty}
              t^q\operatorname{Cap}_S(\{u<-t\};V)=0$;
\item $\displaystyle\lim_{j\to\infty}\nu_{u,j}(V)=0$.
\end{enumerate}
The capacity in (2) can equivalently be replaced by its ordinary-open
outer regularization. If these conditions hold, then
$\mu_{u,j}-\mu_u$ is a finite signed measure on $V$ for every
sufficiently large $j$, and
\[
 \begin{split}
 t^q\operatorname{Cap}_S(\{u<-t\};V)
       &\leq2^q\mu_u(\{u<-t/2\}),\\
 \nu_{u,j}(V)&\leq4^q\mu_u(\{u<-j/4\}),\\
 \|\mu_{u,j}-\mu_u\|_{\mathrm{TV}(V)}
       &\leq(4^q+1)\mu_u(\{u<-j/4\})\longrightarrow0.
 \end{split}
 \tag{BC4}
\]
The measures themselves need not have finite total mass near the
boundary; their difference is supported in one fixed compact set
for all sufficiently large $j$.
\end{theorem}

\noindent\textbf{Status: \texttt{NEW\_PROVED}}. The condition in
(2) is expressed solely through the frozen bounded capacity.

\begin{proof}
Lemma~\ref{lem:r11-compact-truncation-conservation} supplies local
finiteness automatically. For $j>M$, the discarded measure is
carried by $\{u\leq-j\}\subset K_M$. Thus (1) and (3) are
equivalent. Under (1), (UC20) with $v=0$ and $\delta=t/2$
gives the first inequality of (BC4). Its right side tends to zero:
for large $t$ the carrier lies in the fixed compact $K_M$, on
which $\mu_u$ is finite, and the decreasing carriers intersect
in the $\mu_u$-null pole set. Upper semicontinuity makes
$\{u<-t\}$ ordinary open, so its raw and outer capacities agree.
This proves (2).

Conversely, the bounded function $h_j=1+u_j/j$ is a normalized
capacity tester and $\mu_{u,j}=j^q\mu_{h_j}$. Therefore
\[
 \nu_{u,j}(V)
    \leq j^q\operatorname{Cap}_S(\{u<-j/2\};V).
 \tag{BC5}
\]
Condition (2) makes this tend to zero and proves (3).
Combining (BC5) with the first inequality of (BC4), at $t=j/2$,
gives the second. The measures agree on $\{u>-j\}$; on its
complement their variation difference is at most the sum of their
masses. The second inequality and
$\{u\leq-j\}\subset\{u<-j/4\}$ give the last bound in
(BC4). All these tail sets lie in $K_M$ for large $j$, as claimed.
\end{proof}

\begin{theorem}[A solid convex cone of actual potentials]
\label{thm:r11-solid-boundary-cone}
The class $\mathcal F_{S,\partial}(V)$ contains $0$ and has the
following closure properties.
\begin{enumerate}
\item If $u$ belongs to the class and $v$ is any negative extended
classical psh function with $u\leq v\leq0$, then $v$ belongs
to the class.
\item If $u_1,\ldots,u_m$ belong to the class and $a_i\geq0$,
then $\sum_i a_i u_i$ belongs to the class, with zero-coefficient
terms omitted.
\item Finite maxima of its members belong to the class.
\end{enumerate}
No continuity, smoothness, geometric origin or full-support
assumption is imposed on $S$.
\end{theorem}

\noindent\textbf{Status: \texttt{NEW\_PROVED}}. These are
closure assertions for the explicit no-loss class, not for all
unbounded classical psh functions.

\begin{proof}
For (1), $\{v<-t\}\subset\{u<-t\}$ gives both the compact
boundary-gap condition and the capacity-tail criterion.
For (2), assume $A=\sum_i a_i>0$. The elementary inclusion
\[
 \left\{\sum_i a_i u_i<-t\right\}
                  \subset\bigcup_{i:a_i>0}\{u_i<-t/A\}
 \tag{BC6}
\]
gives compact sublevels and
\[
 t^q\operatorname{Cap}_S
          \left(\left\{\sum_i a_i u_i<-t\right\};V\right)
 \leq A^q\sum_{i:a_i>0}
       (t/A)^q\operatorname{Cap}_S(\{u_i<-t/A\};V)
                       \longrightarrow0.
\]
The sum is an extended classical psh function by the frozen
curve definition and positivity. The zero sum is immediate.
Finally a finite maximum is classical psh and dominates any one
of its members, so (1) proves (3).
\end{proof}

\subsection{Top mixed products without a superlinear energy assumption}

\begin{theorem}[A positive mixed top calculus in the boundary class]
\label{thm:r11-boundary-mixed-top-calculus}
For $u_1,\ldots,u_q\in\mathcal F_{S,\partial}(V)$ put
\[
 M_j=T_{S,q}(u_{1,j},\ldots,u_{q,j}),\qquad
 E_j=\bigcap_{i=1}^q\{u_i>-j\}.
\]
The increasing consistent restrictions define a locally finite
positive measure
\[
            M(u_1,\ldots,u_q)=\lim_j\mathbf1_{E_j}M_j.
 \tag{BC7}
\]
It gives zero mass to the union of the $q$ pole sets. The uncut
measures converge to it in total variation on the whole domain,
in the sense of a finite signed difference:
\[
               \|M_j-M(u_1,\ldots,u_q)\|_{\mathrm{TV}(V)}
                                           \longrightarrow0.
 \tag{BC8}
\]
Independent canonical truncation levels have the same limit as
their minimum tends to infinity. The resulting product is symmetric,
additive and positively homogeneous in every slot on this cone, and
\[
 M(u_1,\ldots,u_q)
   ={1\over q!}\sum_{I\subset\{1,\ldots,q\}}
        (-1)^{q-|I|}\mu_{\sum_{i\in I}u_i}.
 \tag{BC9}
\]
Here the empty sum has zero top measure. Every equality is an
equality of actual Borel measures, and hence holds against every
original compact smooth scalar test and all its actual refinements.
\end{theorem}

\noindent\textbf{Status: \texttt{NEW\_PROVED}}. This top-degree
result requires no superlinear weighted energy estimate and does
not assert the corresponding lower-degree recursive calculus.

\begin{proof}
Let $w=\sum_i u_i$ and $s_j=\sum_i u_{i,j}$. The solid-cone
theorem puts $w$ in $\mathcal F_{S,\partial}(V)$. The exact identity
\[
                      \max(s_j,-j)=\max(w,-j)=w_j
 \tag{BC10}
\]
holds: if no slot is truncated then $s_j=w$, while if a slot is
truncated then $s_j\leq-j$. Strict Borel maximum locality thus
identifies the top measure of $s_j$ with $\mu_w$ on $\{w>-j\}$.
As $-qj\leq s_j\leq0$, capacity normalization gives
\[
 \mu_{s_j}(\{w\leq-j\})
       \leq(qj)^q\operatorname{Cap}_S(\{w<-j/2\};V)
                                     \longrightarrow0.
 \tag{BC11}
\]
The mass of $\mu_w$ on the same compact tails tends to zero.
Hence $\mu_{s_j}\to\mu_w$ in total variation on $V$.

The positive bounded mixed expansion gives $q!M_j\leq\mu_{s_j}$.
The complement of $E_j$ is contained in $\{w\leq-j\}$.
Equation (BC11) therefore makes the discarded part
$\mathbf1_{V\setminus E_j}M_j$ tend to zero in mass. Frozen
strict maximum locality gives consistency on $E_j$, exactly as
in (UC2). Their increasing limit is bounded locally by the
already bounded family $\mu_{s_j}/q!$, so is locally finite.
Its mass on $V\setminus E_j$ tends to zero, since these tails
are eventually compact and decrease to the union of the pole sets.
This proves (BC7)--(BC8).

For completeness consider independent levels $l_i\geq m$. Fix
$M>0$ and let $K=\overline{\{w<-M\}}$. Whenever all levels
exceed $M$, all the truncated slots agree with their actual values
outside $K$. A slot-by-slot bounded telescoping identity and
compact Stokes, exactly as in (BC2), show that the integral of
the mixed measure against any fixed cutoff $\chi=1$ near $K$
is independent of all the levels. Passing to the common-level
limit using (BC8) shows that this integral equals $\int\chi\,dM$.
For $m>M$, strict locality gives equality with $M$ on $E_m$,
whereas $V\setminus E_m\subset K$. Subtracting these equal
cutoff integrals yields
\[
 T_{S,q}(u_{1,l_1},\ldots,u_{q,l_q})(V\setminus E_m)
                            =M(V\setminus E_m).
\]
Consequently the total variation of the difference is at most
$2M(V\setminus E_m)$, which tends to zero. This proves independence
of the truncation rates.

Bounded polarization gives (BC9) at level $j$. For every subset
$I$, the preceding argument with $s_{I,j}=\sum_{i\in I}u_{i,j}$
shows $\mu_{s_{I,j}}\to\mu_{\sum_{i\in I}u_i}$ in total
variation; the empty subset contributes zero. Passing in the finite
polarization identity proves (BC9).

Here is also a direct justification of all multilinearity claims.
For any finite family in the cone and fixed positive coefficients
$c_i$, put $A=\sum_i c_i$ and $c=\min_i c_i>0$. With
$s_j=\sum_i c_i u_{i,j}$ and $w=\sum_i c_i u_i$ one has
\[
                    \max(s_j,-cj)=\max(w,-cj).
\]
The preceding locality proof applies, and its discarded mass is
at most $(Aj)^q\operatorname{Cap}_S(\{w<-cj/2\};V)$, which
tends to zero by the tail criterion. Thus the positive bounded
polynomial expansion of $\mu_{s_j}$ passes to the limit for
every set of coefficients. Comparing coefficients in these finite
polynomial identities proves additivity and positive homogeneity
in each slot, with zero coefficients omitted. Symmetry is inherited
from the bounded products. No signed current order or PS quotient
is used.
\end{proof}

\begin{lemma}[Uniform tails above one no-loss lower barrier]
\label{lem:r11-uniform-upper-tail}
If $u,v\in\mathcal F_{S,\partial}(V)$ and $u\leq v\leq0$,
then for every $t>0$,
\[
                 \mu_v(\{v<-t\})
                      \leq2^q\mu_u(\{u<-t/2\}).
 \tag{BC12}
\]
The right side is finite, and tends to zero as $t\to\infty$.
\end{lemma}

\noindent\textbf{Status: \texttt{NEW\_PROVED}} from the
one-sided comparison theorem and the solid boundary cone.

\begin{proof}
Scaling and addition of constants have their expected action on
(UC3): $\mu_{au-b}=a^q\mu_u$ for $a>0$, $b\geq0$.
Indeed, on each $\{u>-k\}$ and for $j>ak+b$, write
$au_k-b=\max((au-b)_j,-ak-b)$ and apply strict maximum
locality and bounded scaling. Exhaust the finite-value locus;
both sides vanish on its complement.
Now $2u$ belongs to the boundary class and
$\mathrm{BG}(2u,v-t)$ holds because its positive-margin gaps lie
in compact strict negative sublevels of $u$. Moreover,
\[
       \{-\infty<v<-t\}\subset\{2u<v-t\}
                                      \subset\{u<-t/2\}.
\]
Apply (UC10) to $2u,v-t$ and use the scaling and translation
identities. The omitted set $\{v=-\infty\}$ is $\mu_v$-null
by its definition; no literal strict-gap inclusion at a common
pole is asserted. This proves (BC12). The limiting assertion follows
from local finiteness on a fixed compact tail and zero mass on
the pole set.
\end{proof}

\begin{theorem}[Joint decreasing continuity in the boundary class]
\label{thm:r11-boundary-joint-decreasing}
For $1\leq i\leq q$, suppose $v_{i,l}$ are negative extended
classical psh functions and
\[
              v_{i,l}\downarrow u_i\in\mathcal F_{S,\partial}(V)
                         \quad\hbox{pointwise on }V.
 \tag{BC13}
\]
Then every $v_{i,l}$ belongs to the same boundary class, and
\[
       M(v_{1,l},\ldots,v_{q,l})
                       \longrightarrow M(u_1,\ldots,u_q)
 \tag{BC14}
\]
against every original compact smooth scalar test. No uniform
superlinear energy bound is needed. The conclusion is ordinary
current convergence, not convergence on arbitrary delta tests.
\end{theorem}

\noindent\textbf{Status: \texttt{NEW\_PROVED}} in this top-degree
no-loss boundary class.

\begin{proof}
Membership follows from upward stability. Put $w=\sum_i u_i$
and $w_l=\sum_i v_{i,l}$, so $w\leq w_l\leq0$. Let
$M_l=M(v_{1,l},\ldots,v_{q,l})$ and let $M_{l,k}$ be the
bounded mixed product with every slot truncated at $-k$.
They agree by strict locality on
$E_{l,k}=\bigcap_i\{v_{i,l}>-k\}$. The positive mixed expansion
and (BC12) give
\[
 M_l(V\setminus E_{l,k})
 \leq {1\over q!}\mu_{w_l}(\{w_l\leq-k\})
 \leq {2^q\over q!}\mu_w(\{w<-k/4\}).
 \tag{BC15}
\]
In the last step use $\{w_l\leq-k\}\subset\{w_l<-k/2\}$
before applying (BC12).
The sum of the $q$ bounded truncated slots lies in $[-qk,0]$.
Also $V\setminus E_{l,k}\subset\{w\leq-k\}$. Capacity
normalization, the positive expansion and the first inequality
of (BC4) therefore give
\[
 M_{l,k}(V\setminus E_{l,k})
 \leq{(qk)^q\over q!}\operatorname{Cap}_S(\{w<-k/2\};V)
 \leq{q^q4^q\over q!}\mu_w(\{w<-k/4\}).
 \tag{BC16}
\]
The right sides are independent of $l$ and tend to zero.
Because the measures agree on $E_{l,k}$, their sum bounds the
total variation of $M_l-M_{l,k}$ on all of $V$.

For fixed $k$, the bounded functions
$\max(v_{i,l},-k)$ decrease jointly to $\max(u_i,-k)$.
Frozen bounded joint decreasing continuity gives convergence of
$M_{l,k}$ against every original compact smooth test as
$l\to\infty$. Apply the uniform tail bounds to the two ends
of this convergence, then let $k\to\infty$ and use (BC8) for
the limiting tuple. This proves (BC14) for the full original
test family. The two limits are separated by explicit uniform
total-variation bounds; no varying-background continuity or
unproved capacity regularity is inserted.
\end{proof}

\begin{corollary}[Outer-capacity domination in the boundary class]
\label{cor:r11-boundary-outer-domination}
Let $u\in\mathcal F_{S,\partial}(V)$ and let $v$ be any
negative extended classical psh function. The strict actual gap
$A=\{u<v\}$ is relatively quasi-open for the fixed capacity:
for each $\eta>0$ there is an ordinary open $G\subset V$ with
$\operatorname{Cap}_S(G;V)<\eta$ such that $A\cup G$ is
ordinary open. In particular,
\[
                 \operatorname{Cap}_S^*(A;V)
                              =\operatorname{Cap}_S(A;V).
 \tag{BC17}
\]
If either $\mu_u(A)=0$ or
$\mathbf1_A\mu_u\leq\mathbf1_A\mu_v$, then
\[
                        \operatorname{Cap}_S^*(A;V)=0.
 \tag{BC18}
\]
If $u,v\in\mathcal F_{S,\partial}(V)$ and $\mu_u=\mu_v$,
then $\operatorname{Cap}_S^*(\{u\ne v\};V)=0$.
These conclusions do not assert actual pointwise equality for a
degenerate background.
\end{corollary}

\noindent\textbf{Status: \texttt{NEW\_PROVED}} from the solid
cone, the new unbounded outer quasicontinuity, and one-sided
domination. Equality of raw and outer capacities is proved only
for these specified strict-gap sets.

\begin{proof}
Put $z=\max(u,v)$. Upward stability makes $z$ a member of the
boundary class even when $v$ itself is not. The exact identity
$A=\{u<z\}$ holds, including the extended values. Since $z\leq0$,
each set $\{u<z-1/m\}$ has compact closure in $V$.
Choose compact neighborhoods $H_m\Subset V$ of these closures,
with ordinary interiors, and set $W=\bigcup_m\operatorname{Int}_X H_m$.
Then $A\subset W$.

Apply Corollary~\ref{cor:r11-unbounded-outer-quasicontinuity} to
$u$ and $z$ on each $H_m$, with total error budget less than
$\eta2^{-m}$. Let $G$ be the union of the resulting ordinary
open exceptional sets. Countable subadditivity gives
$\operatorname{Cap}_S(G;V)<\eta$, and $u,z$ are finite and
continuous on $W\setminus G$: this follows locally from the
cover by the interiors of the $H_m$. Thus $A\setminus G$ is
relatively open in $W\setminus G$. There is an ordinary open
$O\subset W$ with $O\setminus G=A\setminus G$.
Consequently $A\cup G=O\cup G$ is ordinary open. By (UC22),
\[
 \operatorname{Cap}_S^*(A;V)
 \leq\operatorname{Cap}_S(A\cup G;V)
 \leq\operatorname{Cap}_S(A;V)+\operatorname{Cap}_S(G;V).
\]
Let $\eta\downarrow0$ and use the opposite inequality from
monotonicity to prove (BC17).

The condition $\mathrm{BG}(u,v)$ is automatic, since
$\{u<v-\varepsilon\}\subset\{u<-\varepsilon\}$.
Theorem~\ref{thm:r11-unbounded-capacity-domination} therefore
gives raw nullity under either printed measure hypothesis;
(BC17) upgrades it to (BC18). For equal measures and two class
members, apply this argument in both orders and take the union
of the two strict gaps.
\end{proof}

\paragraph{Scope and unsolved gates.}
The boundary class is strictly smaller than all negative extended
classical psh functions: the logarithmic example in
Proposition~\ref{prop:r11-logarithmic-comparison-failure} has
persistent discarded mass and fails the equivalent capacity-tail
criterion. The present top calculus has genuine unbounded members,
for example all radial attenuations $u_a$ with $0<a<1$ from
Section~\ref{sec:r11-radial-unbounded-energy}. It supplies no lower-degree
closedness without additional input, no arbitrary-boundary Perron
solution, and no GQ or CTO relation descent.
